\chapter{Contaminaci\'on Atmosf\'erica}
\section*{Introducci\'on}
\addcontentsline{toc}{section}{Introducci\'on}

La contaminaci\'on atmosf\'erica es uno de los problemas ambientales m\'as cr\'{\i}ticos de la era moderna, con implicaciones significativas para la salud humana, los ecosistemas y el clima global. Este cap\'{\i}tulo aborda las diversas fuentes y tipos de contaminantes que afectan la calidad del aire, as\'{\i} como sus efectos en la salud y el medio ambiente. 

En la Secci\'on~\ref{fgatm}, se exploran las principales fuentes de emisi\'on de gases a la atm\'osfera, clasificadas en biol\'ogicas, la tierra s\'olida, los oc\'eanos y la formaci\'on in situ en la atm\'osfera. Estas fuentes, tanto naturales como antropog\'enicas, contribuyen de manera significativa a la composici\'on qu\'{\i}mica de la atm\'osfera y a la generaci\'on de contaminantes.

La Secci\'on~\ref{sgmog} se enfoca en el fen\'omeno del \textbf{esmog}, un tipo de contaminaci\'on del aire que ha afectado gravemente a regiones como Europa y Los \'Angeles. Adem\'as, se presenta una clasificaci\'on de los contaminantes atmosf\'ericos, destacando su origen y caracter\'{\i}sticas.

En la Secci\'on~\ref{ccriterio}, se analizan los \textbf{contaminantes criterio}, que son sustancias reguladas debido a sus efectos adversos sobre la salud y el medio ambiente. Entre estos se incluyen el mon\'oxido de carbono (\ce{CO}), los compuestos de azufre (\ce{SO2}), los \'oxidos de nitr\'ogeno (\ce{NO_x}), el ozono (\ce{O3}), las part\'{\i}culas suspendidas (\ce{PM10})y el plomo (\ce{Pb}). Para cada uno de estos contaminantes, se discuten sus fuentes de emisi\'on, efectos en la salud y m\'etodos de control.

Finalmente, la Secci\'on~\ref{ctox} aborda los \textbf{compuestos t\'oxicos en el aire}, sustancias que, aunque presentes en concentraciones m\'as bajas, representan un riesgo significativo para la salud humana debido a su alta toxicidad. Se examinan los efectos de estos compuestos y su impacto en la calidad del aire.

Este cap\'{\i}tulo proporciona una visi\'on integral de la contaminaci\'on atmosf\'erica, desde sus fuentes hasta sus consecuencias, con el objetivo de entender mejor este fen\'omeno y promover estrategias para su mitigaci\'on.

\section{Fuentes de gases a la atm\'osfera}
\label{fgatm}
En esta secci\'on se ver\'an las diferentes fuentes que interviene en la emisi\'on de compuestos gaseosos que se encuentran en la atm\'osfera

\subsection{Biol\'ogicas}\index{fuentes!biologicas@biol\'ogicas}\label{subbio}
\begin{itemize}
\item Fotos\'{\i}ntesis de las plantas, responsables del ox\'{\i}geno (\ce{O2}) de la atm\'osfera
\item  Respiraci\'on libera \ce{CO2} en el aire.
\item  Metano (\ce{CH4}), es liberado de materiales org\'anicos, de la fermentaci\'on de vacas, termitas, humedales, cultivos de arroz y tundra.
\item  Terpenos  se evaporan de las hojas, al oxidarse producen part\'{\i}culas y en conjunto con el NOx produce ozono. 
\item Otros gases generados en la biosfera (\ce{CO} y \ce{N2O}) se relacionan con el control de ozono.
\item Quema de biomasa emite gases y part\'{\i}culas a la atm\'osfera. Las especies reactivas de la quema de biomasa inducen un incremento en la producci\'on de ozono y especies del esmog.
\item  Transformaciones biol\'ogicas el \ce{N2} a \ce{NH3}, \ce{N2O} y \ce{NO}.
\item Regiones oce\'anicas de alta productividad biol\'ogica y contenido org\'anico, zonas de surgencia, aguas costeras y las marismas costeras son la principal fuente de bisulfuro de carbono (\ce{CS2})  y sulfuro de carbonilo. (COS). El fitoplancton emite dimetilsulfato (\ce{CH3SCH3} DMS) y bisulfuro de dimetilo (\ce{CH3SSCH3}). La degradaci\'on microbiana de materia org\'anica muerta libera \'acido sulfh\'{\i}drico (\ce{H2S}).
\item Fuentes importantes de cloruro de metilo (\ce{CH3Cl}), que es el halocarbono m\'as abundante en el aire, incluyen la actividad biol\'ogica en el agua de mar, los mohos de la madera y la quema de biomasa. El cloruro de metilo es la principal fuente natural de cloro estratosf\'erico.
\item El hidr\'ogeno molecular se emite por la actividad microbiana en los oc\'eanos, quema de biomasa, la fermentaci\'on bacteriana de suelos y la fotolisis del formaldeh\'{\i}do.
\item El uso por parte de los seres humanos de materiales biog\'enicos da lugar a emisiones de numerosos qu\'{\i}micos t\'oxicos a la atm\'osfera.
\end{itemize}
\begin{ejemplo} {Intercambio de COS} 
Un modelo simple de intercambio de COS entre la capa superficial del oc\'eano, donde se genera, y las regiones inferiores de oc\'eano se pueden basar en las siguientes consideraciones. a) el COS en las capa superficial del oc\'eano est\'a en equilibrio con el COS de la atm\'osfera y no hay sumidero atmosf\'erico neto, b) La concentraci\'on de COS en la superficie del oc\'eano, C(0), es constante a $1.0\times10^{-11}\si{\kilo\gram\per\cubic\metre}$. c) Debajo de la superficie del oc\'eano el COS se destruye qu\'{\i}micamente con  un coeficiente de tasa de primer orden $k=5.0\times10^{-6}\si{\per\second}$. d)  Debajo de la misma capa superficial, la destrucci\'on qu\'{\i}mica de COS se equilibra con el transporte descendente de remolinos con un coeficiente de difusi\'on de \SI{0.0050}{\metre\per\square\second}. e) El oc\'eano es infinitamente profundo y uniforme en la horizontal. Utilice estos supuestos para determinar:
\begin{enumerate}
\item  ?`A qu\'e profundidad la concentraci\'on de COS en el oc\'eano es igual a C(0)/2?
 \item  La densidad de la columna de COS en el oc\'eano.
 \item La vida media de una mol\'ecula de COS en el oc\'eano.
 \end{enumerate}
 
\textbf{Soluci\'on}. Dado que el oc\'eano es uniforme en la horizontal, s\'olo necesitamos considerar la transferencia en la direcci\'on vertical (\textit{z}) (es decir, podemos usar un modelo unidimensional (1D)). Considere una peque\~na porci\'on de oc\'eano a una distancia \textit{z} debajo de la superficie del oc\'eano que tiene una unidad de \'area de secci\'on transversal y espesor $dz$ figura siguiente. Sea $F(z)$ el flujo de COS hacia la parte superior de esta losa y $F(z + dz)$ el flujo que sale de la base de la losa. El flujo neto de COS hacia la losa es:
\begin{equation}
\Delta F=F(z)-F(z+\mathrm{d}z)=F(z)-\left[ F(z)+\frac{\mathrm{d}F(z)}{\mathrm{d}z}\mathrm{d}z \right]=-\mathrm{d}F(z)
\end{equation}
 \end{ejemplo}
\begin{figure}[hbt]
{\centering
\includegraphics[width=0.50\textwidth]{COS_model.png}
\caption{Diagrama del ejercicio}
}
\label{COS_fig}
\end{figure}
\begin{ejemplo} {continuacion}

Como el volumen de la losa es $1\cdot \mathrm{d}z$ 
\begin{equation}
\textrm{el flujo neto de COS en una unidad de volumen de la losa} =-\frac{\mathrm{d}F(z)}{\mathrm{d}z}
\label{flux_COS}
\end{equation}
Si $D$ coeficiente de difusi\'on en remolino,
\begin{equation}
F(z)=-D\frac{\mathrm{d}C(z)}{\mathrm{d}z}
\label{dif_COS}
\end{equation}
donde $C(z)$ es la concentraci\'on de COS a la distancia $z$ debajo de la superficie del oc\'eano, y el signo negativo surge porque la difusi\'on ocurre en la direcci\'on opuesta a la del aumento de la concentraci\'on. En el presente caso, la concentraci\'on de COS es mayor en la superficie del oc\'eano y disminuye hacia abajo; por lo tanto, el flujo de COS es descendente en el oc\'eano. De las ecuaciones. (\ref{flux_COS}) y (\ref{dif_COS})
\begin{equation}
\textrm{el flujo neto de COS en una unidad de volumen de la losa} =D\frac{\mathrm{d}^2C(z)}{\mathrm{d}z^2}
\label{flux2_COS}
\end{equation}
\noindent Si $k$ es el coeficiente de tasa de primer orden para la destrucci\'on qu\'{\i}mica de COS en aguas marinas 

\begin{equation}
\textrm{tasa de destrucci\'on de COS en una unidad de volumen de la losa} =kC(z)
\label{des_COS}
\end{equation}
\noindent En condiciones de estado estacionario, de las ecuaciones (\ref{flux2_COS}) y  (\ref{des_COS})
\begin{equation}
\frac{\mathrm{d}C(z)}{\mathrm{d}t}=D\frac{\mathrm{d}^2C(z)}{\mathrm{d}z^2}-kC(z)=0
\label{cc_COS}
\end{equation}
La soluci\'on de la ecuaci\'on~\ref{cc_COS} es
\begin{equation}
C(z)=C(0)\exp(-z/H)
\label{c_cos}
\end{equation}
\noindent donde H (una altura de escala) est\'a dada por
\begin{equation}
H=\left( \frac{D}{k}\right)^{1/2} =  \left( \frac{0.0050}{5.0\times10^{-6}}\right)=\SI{32}{\metre}
\label{H_COS}
\end{equation}

Ahora se pueden responder las tres preguntas:

\begin{enumerate}
\item Cuando $C(z)=C(0)/2$, tenemos de las ecuaciones (\ref{c_cos}) y  (\ref{H_COS}) 
\begin{equation*}
\frac{1}{2}=\exp\left(-\frac{z}{H} \right)=\exp\left( -\frac{z}{32}\right)
\end{equation*}
Por lo tanto, z = 22m
\item La densidad de la columna de COS en el oc\'eano es la masa total de COS en una columna de unidad de \'area de secci\'on transversal que se extiende desde la superficie hasta el fondo del oc\'eano. Por lo tanto

\begin{eqnarray*}
\textrm{Densidad de columna de COS} &=& \int_{z=0}^\infty C(z)\,\mathrm{d}z \\
       &=&  \int_{z=0}^\infty C(0)\exp\left(-\frac{z}{H}\right)\,\mathrm{d}z  \\
       &=& C(0)H \\
       &=& (1.0\times10^{-11})\times 32 \\
       &=& 3.2\times10^{-10}\si{\kilo\gram\per\square\metre}
\end{eqnarray*}
%
\item De la ecuaci\'on de tiempo de vida de un compuesto $\tau=M/F$
donde $M$ es la cantidad de sustancia qu\'{\i}mica en, digamos, la columna de secci\'on transversal unitaria que consideramos anteriormente y $F$ es la salida (tasa de eliminaci\'on m\'as tasa de destrucci\'on) de la sustancia qu\'{\i}mica de la columna. Por lo tanto, la vida media de una mol\'ecula de COS en el oc\'eano viene dada por:
\begin{eqnarray*}
\tau & = & \frac{C(0)H}{k\int_0^\infty C(z)\,\mathrm{d}z } \\
& =& \frac{C(0)H}{kC(0)H}  \\ 
& =& \frac{1}{k}  \\ 
&=& \frac{1}{5.0\times10^{-6} \si{\second}^{-1}}\\
&=& 2.3\, \textrm{d\'{\i}as} 
\end{eqnarray*}
\end{enumerate}
\end{ejemplo}
\subsection{Tierra s\'olida}\index{fuentes!tierra@tierra s\'olida}
Los volcanes son la fuente geoqu\'{\i}mica m\'as importante de gases traza en al atm\'osfera. Aunque son muy localizados y muy variables, cuando las emisiones volc\'anicas son lanzadas a la estratosfera (o superior), pueden dispersarse r\'apidamente por todo el mundo y tener largos tiempos de residencia (1 a 2 a\~no). El polvo fino  llega a dar varias vueltas alrededor de la Tierra. Los aerosoles producidos a partir de los gases pueden permanecer en la estratosfera durante varios a\~nos. Esto, a su vez, da lugar a un enfriamiento medio global en la superficie de $\sim$ \SI{0.4}{\celsius} en el a\~no siguiente a la erupci\'on; este enfriamiento llega a desaparecer despu\'es de $\sim$3 a\~nos a medida que disminuye el polvo en la atm\'osfera. Adem\'as de cenizas y part\'{\i}culas de polvo, los volcanes emiten \ce{H2O}, \ce{CO2}, \ce{SO2}, \ce{H2S}, COS, \ce{CS2}, HCl, HF, HBr, \ce{CH4}, \ce{CH3Cl}, \ce{H2}, CO y varios metales pesados vol\'atiles (p. ej., Hg). Las rocas son una fuente de  cantidades peque\~nas de ciertos gases y son las principales fuentes de los gases He, Ar y Rn en la atm\'osfera. El helio se produce principalmente del decaimiento radiactivo del uranio-238 y torio-232- No se acumula en la atm\'osfera debido a que es tan ligero que escapa de la exosfera\index{exosfera}. El arg\'on se ha acumulado durante eones debido a la desintegraci\'on radiactiva del potasio-40 de rocas.

Las rocas carbonatadas, como la piedra caliza, que se presenta principalmente como calcita (es decir, \ce{CaCO3}), contienen 100,000 veces m\'as carbono que la atm\'osfera, pero la mayor parte est\'a secuestrada. Sin embargo, las rocas carbonatadas y los sedimentos marinos participan en un ciclo de per\'{\i}odo largo con el \ce{CO2} atmosf\'erico como se muestra a continuaci\'on. La erosi\'on de la calcita por el \ce{CO2} disuelto en la lluvia y en aguas dulces (r\'{\i}os y lagos) se puede representar por

\begin{footnotesize}
\begin{eqnarray}
\ce{CO2}(g) + \ce{H2O} &\rightleftarrows& \ce{CO2}(ac) + \ce{H2O}(l)  \label{eq_co2h2o}\\
\ce{2H2O}(l) +   \ce{CO2}(ac)&\rightleftarrows&   \ce{H3O^+}(ac) +  \ce{HCO3^-}(ac) \label{eq_co2}\\
\ce{CaCO3}(s) +   \ce{H3O^+}(ac) +  \ce{HCO3^-}(ac) &\rightleftarrows& \ce{Ca^{2+}}(ac) + \ce{2HCO3^-}(aq) + \ce{H2O}(l)\label{eq_caco3}\\\hline
\textrm{Neto:       }   \ce{CaCO3}(s) +\ce{CO2}(g) + \ce{H2O}(l)  &\rightleftarrows&\ce{Ca^{2+}}(ac) + \ce{2HCO3^-}(aq)  \label{eq_cca}
\end{eqnarray}
\end{footnotesize}

La reacci\'on (\ref{eq_co2h2o}) representa el equilibrio entre el \ce{CO2} en el aire y en los r\'{\i}os y lagos de agua dulce. En la reacci\'on (\ref{eq_co2}), el \ce{CO2} recibe un \ce{OH^-}  del \ce{H2O} para formar el \'acido muy d\'ebil (es decir, poco ionizado) \ce{HCO3^-} (el ion bicarbonato), y en la reacci\'on (\ref{eq_caco3}) el \ce{H3O+} remueve \ce{CO3^{2-}} del \ce{CaCO3} para formar otro ion bicarbonato\index{bicarbonato}. La reacci\'on directa de (\ref{eq_caco3}) representa la erosi\'on de la calcita por el \ce{CO2}  disuelto en aguas dulces.

Los productos de la \gls{meteorizacion} ,en el lado derecho de la reacci\'on~\ref{eq_cca}, eventualmente ingresan a los oc\'eanos, donde se precipitan para formar nuevos sedimentos (lo contrario de la reacci\'on (\ref{eq_cca})). A trav\'es del levantamiento de las regiones de la plataforma continental, la subducci\'on de sedimentos marinos hacia el manto superior y la corteza inferior de la Tierra y las erupciones volc\'anicas, estos productos eventualmente regresan a los sedimentos continentales, completando as\'{\i} este ciclo geoqu\'{\i}mico. Los tiempos de residencia del \ce{Ca^{2+}}(aq) y del \ce{HCO3^-}(aq) en los oc\'eanos son $\sim 8\times 10^5$ y $\sim7.5 \times 10^4$ a\~nos, respectivamente.

\begin{ejemplo}{Disoluci\'on de \ce{CO2} en agua} 
Cuando el \ce{CO2} se disuelve en agua pura \textquestiondown qu\'e reacciones se producen?

\textbf{Respuesta:}

\begin{footnotesize}
\begin{eqnarray}
\ce{CO2}(g) + \ce{H2O} &\rightleftarrows& \ce{H2CO3}(ac) \\
\ce{H2CO3}(ac) + \ce{H2O}(l) &\rightleftarrows&  \ce{HCO3^-}(ac) + \ce{H3O^+}(ac) \\
 \ce{HCO3^-}(ac) + \ce{H2O}(l) &\rightleftarrows&  \ce{CO3^{-2}}(ac) + \ce{H3O^+}(ac) \\ \hline
 \textrm{Neto:         } \ce{CO2}(g) + 3\ce{H2O} &\rightleftarrows& \ce{CO3^{-2}}(ac) +2 \ce{H3O^+}(ac) 
 \end{eqnarray}
\end{footnotesize}
\end{ejemplo}

\subsection{Oc\'eanos}\index{fuentes!oceano@oc\'eano}
La mayor\'{\i}a de los gases que pasan de los oc\'eanos al aire se originan en procesos biol\'ogicos, que se analizan en la subsecci\'on (\ref{subbio}). Adem\'as, como vimos en la subsecci\'on anterior, los oc\'eanos pueden participar en el ciclo de gases entre la Tierra s\'olida y la atm\'osfera. 

Los oc\'eanos son un enorme reservorio de aquellos gases que son apreciablemente solubles en agua. Los oc\'eanos pueden servir como fuente y sumidero de gases traza en la atm\'osfera.

\subsection{Formaci\'on in situ en la atm\'osfera.}\index{atmosfera@atm\'osfera!fuentes de contaminantes}

Las reacciones qu\'{\i}micas en la atm\'osfera son una fuente importante y un sumidero de muchos constituyentes traza. Los gases traza emitidos por la biosfera, la Tierra s\'olida y los oc\'eanos generalmente se encuentran en un estado de oxidaci\'on reducido (es decir, m\'as bajo) (por ejemplo, C, N y S), pero cuando regresan a la superficie de la Tierra generalmente se encuentran en un estado de  mayor oxidaci\'on.

Los procesos que producen estas transformaciones se pueden dividir en reacciones de fase gaseosa homog\'enea, fase acuosa homog\'enea y reacciones heterog\'eneas.

 \section{El esmog}\index{esmog}
\label{sgmog}

 La contaminaci\'on ambiental se encuentra presente en nuestra vida y es producto de las actividades de la vida cotidiana tanto en las ciudades como en el campo. La principal causa de la contaminaci\'on en el aire es debido a la combusti\'on, y \'esta es de suma importancia para el ser humano. Cuando ocurre la combusti\'on perfecta o te\'orica, el carbono e hidr\'ogeno del combustible se combinan con el ox\'{\i}geno del aire desprendiendo calor, luz y generando bi\'oxido de carbono y vapor de agua. Sin embargo las impurezas del combustible, el mezclado imperfecto, la temperatura exterior afectan la combusti\'on gener\'andose impurezas tales como mon\'oxido de carbono, di\'oxido de azufre, \'oxidos de nitr\'ogeno, holl\'{\i}n y combustible no quemado ---todos ellos son contaminantes del aire.
 
 El esmog es la combinaci\'on de part\'{\i}culas de humo de fuentes industriales con neblina produciendo un color negro-amarillento cerca de la superficie (nieblumo)\index{esmog!nieblumo}. La combusti\'on de gasolina puede crear otro problema de contaminaci\'on conocido como ``esmog fotoqu\'{\i}mico''. Este se genera cuando los contaminantes primarios (\'oxidos de nitr\'ogeno y COV) interact\'uan bajo la influencia de la luz solar para producir una mezcla de cientos de sustancias t\'oxicas diferentes conocidas como contaminantes secundarios. Es una bruma marr\'on que contamina las ciudades. Puede hacer dif\'{\i}cil el respirar para algunas personas y reduce la visibilidad.

 \subsection{Europa}
En diciembre de 1930, una regi\'on altamente industrializada del Valle de Meuse, en B\'elgica, se cubri\'o durante 3 d\'{\i}as de una espesa niebla, por lo que cientos de personas enfermaron y 60 murieron ---m\'as de 10 veces el n\'umero normal.  

En Inglaterra durante 9 d\'{\i}as, en enero de 1931, murieron 592 personas --lo que nuevamente representaba un considerable incremento en la tasa de mortalidad. En 1948, en Donora, Pennsylvania, un peque\~no pueblo en donde hab\'{\i}a plantas qu\'{\i}micas y acerer\'{\i}as se cubri\'o por una niebla durante 4 d\'{\i}as,  enferm\'o casi la mitad de sus 14,000 habitantes. Murieron 20 personas. No fue hasta que una gran capa de niebla cubri\'o Londres en 1952 cuando se hizo totalmente evidente el siniestro potencial de la contaminaci\'on del aire. La niebla dur\'o desde el 5 de diciembre hasta el 8, y 10 d\'{\i}as despu\'es se supo que el n\'umero total de muertes en la regi\'on principal de Londres sobrepasaba en 4,000 al promedio. Las estad\'{\i}sticas indicaron que casi todos los que hab\'{\i}an muerto inesperadamente ten\'{\i}an antecedentes cl\'{\i}nicos de bronquitis, enfisema o trastornos card\'{\i}acos, y que las personas clasificadas en la \'ultima categor\'{\i}a, eran las m\'as vulnerables. Nuevamente, en enero de 1956, se produjeron 1,000 muertes m\'as debido a una extensa niebla.

 
 \subsection{Los Angeles}
En contraste con la contaminaci\'on en Londres los episodios de contaminaci\'on en Los Angeles y la Ciudad de M\'exico se dan en situaciones de cielos despejados con sistemas de alta presi\'on que limitan la dispersi\'on de contaminantes. En este caso las emisiones primarias de los procesos industriales, comerciales y del transporte reaccionan en la atm\'osfera generando nuevos contaminantes que afectan la salud de los seres humanos, animales, plantas  y afectando a los materiales.

\subsection{Clasificaci\'on de Contaminantes}

Los contaminantes atmosf\'ericos se pueden clasificar de diversas maneras, una de ellas es por su composici\'on:\index{contaminante!clasificaci\'on}

\begin{itemize}
\item \textbf{Org\'anicos} Contienen carbono, hidr\'ogeno (\ce{CH4}) y pueden incluir ox\'{\i}geno (HCHO), nitr\'ogeno (\ce{CH3NH2}), azufre (DMS).
\item \textbf{Inorg\'anicos} Compuestos de azufre (\ce{SO2}, \ce{H2S}), nitr\'ogeno (\ce{NO_$x$}), halogenados (CFC)\index{CFC}, metales pesados (As, Pb, Cd, etc)
\end{itemize}
Por su estado f\'{\i}sico:
\begin{itemize}
\item Gases y vapores
\item Part\'{\i}culas (s\'olidas y l\'{\i}quidas).
\end{itemize}

A su vez las part\'{\i}culas se pueden clasificar de acuerdo a con su tama\~no, como se muestra en el \textbf{Cuadro~\ref{part:1}} en p\'agina~\pageref{part:1}

Tambi\'en se puede categorizar los contaminantes por su origen de la siguiente  forma:\index{contaminante!primarios}\index{contaminante!secundario}
\begin{description}
\item[Primarios] Aquellos que son emitidos directamente
dentro de la atm\'osfera. Como lo son el mon\'oxido de nitr\'ogeno (NO), di\'oxido de azufre (\ce{ SO2}), el mon\'oxido de carbono (CO).
\item[Secundarios] Creados a trav\'es de reacciones y procesos en la
atm\'osfera. Como el ozono (\ce{O3}), tri\'oxido de azufre (\ce{SO3}) y el nitrato de
peroxiacilo (PAN)\index{PAN}.
\end{description}
\section{Contaminantes criterio}\index{contaminante!criterio}
\label{ccriterio}

De todos los contaminantes posibles existe un conjunto de ellos que se emplean para calificar la calidad del aire  y se les conoce como  \textit{contaminantes criterio}, los contaminantes criterio son seleccionados debido a que pueden provocar efectos a la salud y son los que se muestran en el  \textbf{Cuadro~\ref{contcrit}}.
\begin{table}[htb]
\caption{Contaminantes Criterio}
\label{contcrit}
\begin{center}
\begin{tabular}{ll}\hline
Mon\'oxido de carbono & CO\\
Di\'oxido de azufre & \ce{SO2}\\
\'oxidos de nitr\'ogeno &NO$_x$\\
Ozono & \ce{O3}\\
Part\'{\i}culas & \gls{PST}/\ce{PM10},\ce{PM_{2.5}}\\
Plomo& Pb\\ \hline
\end{tabular}
\end{center}
\end{table}

De todos los compuestos que se han presentado los hidrocarburos (COV)\index{COV} no son contaminantes criterio, pero se estudian muy de cerca por su rol en la producci\'on de ozono:

\begin{center}
COVs $+$ NO$_x +$luz solar $\longrightarrow$ esmog fotoqu\'{\i}mico (\ce{O3}) 
\end{center}

\subsection{Mon\'oxido de carbono (CO)}
\index{CO}
Es gaseoso, inodoro e incoloro es resultado de la combusti\'on incompleta de los veh\'{\i}culos, el 70\% del CO proviene de los veh\'{\i}culos y presenta un perfil diurno, con concentraciones m\'aximas en horas pico de tr\'ansito.

\subsubsection{Efectos a la salud:}
\index{CO}
El CO se puede ligar a la sangre en lugar del di\'oxido de carbono (\ce{CO2}) por lo que puede ser asfixiante al reducir el ox\'{\i}geno (\ce{O2}) transportado -- disminuye la actividad cerebral, incrementa la frecuencia card\'{\i}aca. Los efectos dependen del nivel de actividad.
\begin{equation}
\%\ce{COHb}  = 0.005[\ce{CO]}^{0.85}(\alpha t)^{0.63}
\end{equation}
Donde:

\begin{tabular}{r @{---}l}
$\%\gls{COHb} $ & Carboxihemoglobina en la sangre en por ciento de sa\-tu\-ra\-ci\'on.\\
$[\ce{CO}]$ &Concentraci\'on de CO en ppm\\
$\alpha$ & Coeficiente de actividad ( 0.94 sedentario, 3 trabajo pesado)\\
$t$ & tiempo de exposici\'on en minutos\\
\end{tabular}

Algunos efectos por la cantidad de CO en la sangre son los siguientes:
\begin{itemize}
\item 2.5\% se muestran algunos efectos ligeros.
\item 5\% Se notan efectos psicomotores.
\item 10\% Mareos y dolor de cabeza.
\item 50\% Mortal.
\end{itemize}

\subsubsection{Fuentes de exposici\'on:}
Intersecciones muy transitadas, vivir cerca de vialidades transitadas, fumar.
Se requieren de 3 a 4 horas para alcanzar niveles altos de concentraci\'on de COHb, y
se requieren de 3 a 4 horas para recobrarse de la exposici\'on.

\subsection{Compuestos de azufre}\index{azufre!compuestos}
Este contaminante puede producir, incluso a grandes distancias de la fuente de emisi\'on, efectos adversos sobre la salud (tales como irritaci\'on e inflamaci\'on del sistema respiratorio, afecciones e insuficiencias pulmonares, alteraci\'on del metabolismo de las prote\'{\i}nas, dolor de cabeza o ansiedad), sobre la biodiversidad, los suelos y los ecosistemas acu\'aticos y forestales (puede ocasionar da\~nos a la vegetaci\'on, degradaci\'on de la clorofila, reducci\'on de la fotos\'{\i}ntesis y la consiguiente p\'erdida de especies) e incluso sobre las edificaciones, a trav\'es de procesos de acidificaci\'on, pues una vez emitido, reacciona con el vapor de agua y con otros elementos presentes en la atm\'osfera, de modo que su oxidaci\'on en el aire da lugar a la formaci\'on de \'acido sulf\'urico.

Adem\'as, tambi\'en act\'ua como precursor de la formaci\'on de sulfato am\'onico, lo que incrementa los niveles de \ce{PM10} y \ce{PM_{2.5}}, con graves consecuencias igualmente sobre la salud.

\subsubsection{Fuentes de emisi\'on del di\'oxido de azufre} Las principales fuentes de azufre en la atm\'osfera son el uso de combustibles f\'osiles y las exhalaciones volc\'anicas\index{volcan@volc\'an}. El compuesto que principalmente se emite es el di\'oxido de azufre (\ce{SO2}). La fuente  \gls{antropica} principal es la quema de combustibles conteniendo azufre, el 85\% del total de emisiones proviene de esta fuente. Algunas industrias de la fundici\'on, la refinaci\'on, la minera --  tienen normas de emisi\'on de este compuesto.

El contenido de azufre en los combustibles abarca el intervalo de 0.05 a 6\% en los combustibles l\'{\i}quidos y en el carb\'on. Cuando se queman se emite el di\'oxido de azufre (\ce{SO2}).
Produce lluvia \'acida\index{lluiva@lluvia \'acida}, tambi\'en se condensa en part\'{\i}culas formando \ce{SO_4^=} durante el per\'{\i}odo de d\'{\i}as
\begin{itemize}
\item Exposici\'on a 1 ppm de \ce{SO2} por 15--20 min se observan cambios en el pulso y
respiraci\'on.
\item 0.3 ppm por 1--3 d\'{\i}as enfermedades cardio-respiratorias, muertes en Londres.\index{Londres}
\item 0.002 ppm por un a\~no, incrementa las admisiones al hospital.
\end{itemize}

\subsection{Compuestos de nitr\'ogeno}
\index{nitrogeno@nitr\'ogeno!compuestos de} 
En la atm\'osfera se encuentran diversas formas del nitr\'ogeno, tanto oxidadas como reducidas:

\textbf{Formas oxidadas:}
\begin{itemize}
\item NOx= NO + \ce{NO2},
\item NOy =NOx + \ce{NO3} + HONO + \ce{HNO3} + \ce{N2O5} + PAN
\item NOz = NOy -NOx
\end{itemize}

\textbf{Formas reducidas:}
\begin{itemize}
\item  Amon\'{\i}aco\index{amoniaco@amon\'{\i}aco} (\ce{NH3})
\item Cianuro de hidr\'ogeno (HCN)\index{cianuro} 
\item Aminas:
\begin{itemize}
\item Alif\'aticas\index{aminas!alifaticas@alif\'aticas}: \ce{RNH2} 
\item Arom\'aticas\index{aminas!aromaticas@arom\'aticas}  RR'NH y RR'R''N 
\end{itemize}
\item N\'{\i}tratos: RCN
\end{itemize}
 donde R, R' R'' = son grupos alquilos\index{grupo!alquilo}  o \gls{arilo}\index{grupo!arilo} .

\subsubsection{\'oxidos de Nitr\'ogeno (NO$_x$)}
La suma de \ce{NO} + \ce{NO2} se refiere como NO$_x$, principalmente proviene de la combusti\'on:
\begin{description}
\item[NO$_x$ Termal] Cuando el \ce{N2}  y el \ce{O2} se calientan a altas temperaturas
($>400^\circ C$)
\item[NO$_x$ Combustible] Del nitr\'ogeno contenido en el combustible.
\end{description}
La emisiones de NO$_x$ son principalmente NO, \'este reacciona y forma el \ce{NO2} que absorbe la luz.
%\begin{eqnarray}
%\ce{NO2} + \textrm{hidrocarburos}&\longrightarrow&\textrm{esmog}\index{esmog}\\
%\ce{NO2}+ \ce{OH}\cdot &\longrightarrow&\ce{HNO3} \: \textrm{(Lluvia \'acida)} \index{lluiva@lluvia \'acida}
%\end{eqnarray} 

\subsubsection{Control} En emisiones de fuentes estacionarias y en motores de veh\'{\i}culos. Reduciendo la temperatura de combusti\'on y con combustibles con bajo nitr\'ogeno (N). Altas temperaturas disminuyen el mon\'oxido de carbono (CO) pero incrementan la emisi\'on de  NO$_x$.

\subsection{Ozono (O$_3$)}\label{o3_salud}\index{ozono}
El ozono no se emite directamente a la atm\'osfera, se produce mediante una serie de reacciones que involucran a los \'oxidos de nitr\'ogeno, hidrocarburos y la luz solar ($h\nu$). El ozono posee dos efectos -- da\~nos a la salud y a los materiales.

Efectos del ozono en materiales:
\begin{itemize}
\item Reduce la vida \'util de llantas y hules.
\item Puede da\~nar a la vegetaci\'on (reduce la producci\'on agr\'{\i}cola)
\end{itemize}
Efectos del ozono en la salud:
\begin{itemize}
\item Irritaci\'on de ojos.
\item Constricci\'on del pecho, irritaci\'on en la garganta.
\item Altas concentraciones agravan las enfermedades respiratorias.
\end{itemize}
La  formaci\'on de ozono depende de la disponibilidad del \ce{NO2}  y la destrucci\'on depende de la concentraci\'on del NO. Sin embargo en las ciudades el efecto global de las reacciones con hidrocarburos (COV) es convertir el \ce{NO} a \ce{NO2} por lo cual hace que se incremente la concentraci\'on del \ce{O3}.

\begin{table}[hbt]
\caption{Ozono en la atm\'osfera}
\label{ozon:1}
\begin{center}
{\footnotesize \begin{tabularx}{.9\textwidth}{X|X}
\textbf{Ozono Troposf\'erico}&\textbf{Ozono Estratosf\'erico}\\ \hline \hline
Localizado entre $0$--$10$ \si{\kilo\metre}& Localizado entre $15$--$35$ \si{\kilo\metre}\\
Contiene el $10$\% del ozono de la atm\'osfera&
Contiene el $90$\% del ozono de la atm\'osfera\\
Impacto da\~nino: efectos t\'oxicos en seres humanos y vegetaci\'on&
Rol ben\'efico: Act\'ua como escudo a la radiaci\'on UV\\
T\'opicos actuales: -- Episodios de alta concentraci\'on en superficie en ciudades.&
T\'opicos actuales: Disminuci\'on en la concentraci\'on. Hoyo de ozono en la
primavera ant\'artica\\\hline
\end{tabularx}}
\end{center}
\end{table}

El control del ozono se da a trav\'es de controlar sus precursores, \'oxidos de nitr\'ogeno, y compuesto org\'anicos. Si el aire se encuentra ``atrapado'' por la
orograf\'{\i}a y condiciones meteorol\'ogicas como la inversi\'on t\'ermica, esto puede empeorar el problema como en  Los Angeles, la Ciudad de M\'exico, Grecia y en el Valle Fraser (frontera EE-Canad\'a).

\paragraph{Nota sobre ozono}
\index{ozono!troposf\'erico}
\index{ozono!estratosf\'erico}
Existen dos diferentes zonas donde se encuentra el ozono, el troposf\'erico (nivel del suelo) y el estratosf\'erico. En el \textbf{Cuadro~\ref{ozon:1}} se pueden observar las
caracter\'{\i}sticas y efectos que posee el ozono en las diferentes zonas.



\subsection{Part\'{\i}culas}
Se definen como part\'{\i}culas a cualquier material disperso (s\'olido o l\'{\i}quido) cuyos componentes individuales tienen un tama\~no mayor que una mol\'e\-cula pero menor a \SI{500}{\micro\metre}.  En el \textbf{Cuadro \ref{part:1}}  se enlistan en diferentes categor\'{\i}as las part\'{\i}culas:

\begin{table}[h!]
\caption{Definici\'on de las part\'{\i}culas suspendidas en el aire}
\label{part:1}
\begin{center}
{\footnotesize \begin{tabularx}{\linewidth}{>{\setlength{\hsize}{.35\hsize}}X%
>{\setlength{\hsize}{1.65\hsize}}X} \hline
{\footnotesize Part\'{\i}culas} & Cualquier material, excepto agua no combinada, que
existe en estado s\'olido o l\'{\i}quido en la atm\'osfera o en una corriente de gas
en condiciones normales.\\
{\footnotesize Aerosol} & Una dispersi\'on de part\'{\i}culas microsc\'opicas, s\'olidas o
l\'{\i}quidas. \\
{\footnotesize  Polvo} & Part\'{\i}culas s\'olidas de un tama\~no mayor que el coloidal, capaces
 de estar en suspensi\'on temporal en el aire.\\
{\footnotesize Ceniza fin}a& Part\'{\i}culas de ceniza finamente divididas arrastradas
por el gas de combusti\'on. Las part\'{\i}culas pueden contener combustible no
quemado\\ 
{\footnotesize  Niebla} &Aerosol visible.\\
{\footnotesize Vapores}& Part\'{\i}culas formadas por la condensaci\'on, sublimaci\'on,
o reacci\'on qu\'{\i}mica, predominantemente mayores de \SI{1}{\micro\metre} (humo de
tabaco).\\
{\footnotesize Neblina} & Dispersi\'on de peque\~nas gotas de l\'{\i}quido de suficiente tama\~no como
para caer.\\
 {\footnotesize Part\'{\i}cula}& Masa discreta de materia s\'olida o l\'{\i}quida.\\
{\footnotesize Humo}& Part\'{\i}culas peque\~nas arrastradas por los gases, que resultan
dela combusti\'on.\\
{\footnotesize Holl\'{\i}n}& Una aglomeraci\'on de part\'{\i}culas de carb\'on.\\ \hline
\end{tabularx}}
\end{center}
\end{table}
\subsubsection{Efectos en la salud} 
Las part\'{\i}culas menores a \SI{10}{\micro\metre} (PM$_{10}$) y \SI{2.5}{\micro\metre}  (PM$_{2.5}$) constituyen el principal riesgo para la salud humana. Al ser inhaladas, penetran en el tracto respiratorio: las PM$_{10}$ se depositan en v\'{\i}as superiores, mientras las PM$_{2.5}$ alcanzan los alv\'eolos pulmonares. Esto desencadena procesos inflamatorios sist\'emicos, reducci\'on de la capacidad pulmonar (2-3\% por cada \SI{10}{\micro\gram\per\cubic\metre}  de PM$_{2.5}$) y exacerbaci\'on de enfermedades cardiovasculares. La exposici\'on cr\'onica se relaciona con:
\begin{itemize}
    \item Incremento del 12\% en mortalidad por c\'ancer de pulm\'on (por cada \SI{10}{\micro\gram\per\cubic\metre} anual de PM$_{2.5}$) \index{mortalidad!particulas@part\'{\i}culas}
    \item Desarrollo de EPOC\footnote{Enfermedad pulmonar obstructiva cr\'onica} en poblaciones urbanas (OR = 1.29; IC95\%: 1.04–1.59)\index{EPOC}
    \item Neurotoxicidad en ni\~nos (p\'erdida de 1-2 puntos de CI por cada \SI{5}{\micro\gram\per\cubic\metre} de PM$_{2.5}$)
\end{itemize}

\paragraph{Nota sobre las part\'{\i}culas}
Cuando se especific\'o la norma de part\'{\i}culas, \'estas se regular\'on como Part\'{\i}culas Suspendidas Totales (PST)\index{PST}. Esta fue la medida de concentraci\'on de part\'{\i}culas en el aire.\index{particulas@part\'{\i}culas}

En 1987 en EU la norma cambi\'o a PM$_{10}$ -- toda part\'{\i}cula menor a  \SI{10}{\micro\metre} en di\'ametro y en 1997, otra norma de part\'{\i}culas se adicion\'o --- \ce{PM_{2.5}} -- toda part\'{\i}cula menor a \SI{2.5}{\micro\metre}  en di\'ametro.

\subsection{Plomo (Pb)}
\index{plomo}
Hasta 1980, la principal fuente de emisiones de plomo atmosf\'erico proven\'{\i}a de los veh\'{\i}culos que utilizaban \textbf{gasolina con tetraetilo de plomo} (\ce{(C2H5)4Pb}), un aditivo antidetonante. \index{antidetonante}Este compuesto se a\~nad\'{\i}a en concentraciones de \textbf{1.1 gramos por gal\'on} (\SI{ 0.29}{\gram\per\litre}), equivalente al 61\% en peso del aditivo completo \textit{TEL} (Tetraethyl Lead). \index{tetraetilo de plomo}

\subsubsection{Efectos en la salud} El plomo representa un riesgo toxicol\'{o}gico significativo, particularmente en poblaciones infantiles. Sus principales v\'{\i}as de exposici\'{o}n incluyen:

\begin{itemize}
    \item \textbf{Inhalaci\'{o}n}: 
    \begin{itemize}
        \item Emitido como sales de plomo (ej. \ce{PbSO4}) depositadas cerca de v\'{\i}as de tr\'{a}nsito
        \item Part\'{\i}culas transportadas a interiores por calzado (tama\~no $\sim$ \SIrange{2}{10}{\micro\metre})
        \item Resuspensi\'{o}n atmosf\'{e}rica en espacios cerrados (factor $3--5\times$ mayor que exterior)
    \end{itemize}
    
    \item \textbf{Ingesti\'{o}n}:
    \begin{itemize}
        \item Contaminaci\'{o}n h\'{\i}drica por tuber\'{\i}as de plomo (pH $<$ 7 aumenta lixiviaci\'{o}n)
        \item Pinturas con base de plomo (comunes en construcciones anteriores a 1980)
        \item Suelos contaminados (bioaccesibilidad $\sim$ 25-60\%)
    \end{itemize}
\end{itemize}

\textbf{Mecanismo de toxicidad}:
\begin{itemize}
    \item Sustituci\'{o}n del hierro en rutas metab\'{o}licas (afinidad por grupos sulfhidrilo)
    \item Neurotoxicidad aguda en ni\~nos (DL$_{50}$ = \SIrange{100}{150}{\micro\gram\per\deci\liter}):
    \begin{itemize}
        \item Reducci\'{o}n de coeficiente intelectual (2-5 puntos por \SI{10}{\micro\gram\per\deci\liter})
        \item Da\~no cerebral irreversible (mielinizaci\'{o}n alterada)
        \item Trastornos conductuales (TDAH\footnote{Trastorno de deficit de atenci\'on e hiperactividad.}, agresividad)
    \end{itemize}
    
    \item Umbrales cr\'{\i}ticos (CDC 2023):
    \begin{itemize}
        \item \SI{20}{\micro\gram\per\deci\liter}: Da\~nos hematol\'{o}gicos (anemia hemol\'{\i}tica)
        \item \SI{50}{\micro\gram\per\deci\liter}: Encefalopat\'{\i}a aguda (mortalidad $\sim$ 5-10\%)
    \end{itemize}
\end{itemize}

\subsection{Visibilidad y la calidad del aire}\index{visibilidad!calidad del aire}

La dispersi\'on de la radiaci\'on electromagn\'etica en la atm\'osfera est\'a gobernada por el par\'ametro de tama\~no $\alpha = \frac{2\pi r}{\lambda}$, donde $r$ es el radio de la part\'{\i}cula y $\lambda$ la longitud de onda. Cuando $\alpha \approx 1$ (dispersi\'on de Mie), se produce la m\'axima eficiencia de dispersi\'on. Para la luz visible ($\lambda = \SIrange{0.38}{0.70}{\micro\metre}$), esto corresponde a part\'{\i}culas con di\'ametros d = \SIrange{0.12}{0.22}{\micro\metre}.
\subsubsection{Par\'ametros operacionales}\index{visibilidad!horizontal}
\begin{itemize}
    \item \textbf{Visibilidad horizontal} (f\'ormula de Koschmieder\index{Koschmieder}):\label{fKoschmieder}
    \begin{equation}
    V = \frac{3.912}{\beta_{ext}} \quad [\si{\kilo\metre}], \quad \beta_{ext} = \sum N_i\sigma_{ext,i}
    \end{equation}
    
    \item \textbf{Umbrales significativos}:
    \begin{itemize}
        \item $2,000$ part/cm$^3$: Obscurecimiento de relieves monta\~nosos
        \item $100,000$ part/cm$^3$: $V \approx \SI{1.6}{\kilo\metre}$ (1 milla).
        \item $8-10$ ppm de \ce{NO2}: $V \leq \SI{1.6}{\kilo\metre}$.
    \end{itemize}
    
    \item \textbf{Efecto de humedad}:
    \begin{itemize}
        \item HR > 70\%: Crecimiento higrosc\'opico ($\times 1.5-4$ en $d$)
        \item Aumento equivalente en $\beta_{ext}$ de 2--5 veces
    \end{itemize}
\end{itemize}

\subsubsection{Contexto regulatorio}\footnote{
Evoluci\'on de normas de part\'{\i}culas:
\begin{itemize}
    \item PST (Part\'{\i}culas Suspendidas Totales): Medida hist\'orica.
    \item PM$_{10}$ (1987 EU): Part\'{\i}culas < \SI{10}{\micro\metre} .
    \item PM$_{2.5}$ (1997 EU): Part\'{\i}culas < \SI{2.5}{\micro\metre}.
\end{itemize}
}

Los principales gases que influyen en la visibilidad son \ce{SO2}, \ce{NO}, y \ce{H2O}, donde:
\begin{itemize}
    \item \ce{SO2}: Contribuye a formaci\'on de aerosoles sulfatados
    \item \ce{NO}: Precursor de \ce{NO2} y nitratos
    \item \ce{H2O}: Modula crecimiento de part\'{\i}culas higrosc\'opicas
\end{itemize}

En condiciones pr\'{\i}stinas ($\beta_{ext} < 14\ \text{Mm}^{-1}$) se alcanzan $V > \SI{280}{\kilo\metre}$, mientras que en niebla densa ($\beta_{ext} > 39,000\ \text{Mm}^{-1}$) $V < \SI{0.1}{\kilo\metre}$.

\section{Compuestos t\'oxicos en el aire}
\label{ctox}
\normalsize
El t\'ermino de \emph{compuestos t\'oxicos} se refiere a cualquier sustancia qu\'{\i}mica presente en el aire con potencial de da\~no a la salud o el ambiente. Para prop\'ositos regulatorios, la Agencia de Protecci\'on Ambiental de EE.UU. (EPA) clasifica estos contaminantes en dos categor\'{\i}as:

\begin{itemize}
    \item \textbf{Contaminantes criterio}: 
    \begin{itemize}
        \item Regulados mediante Normas Nacionales de Calidad del Aire (NAAQS) en EEUU y por Normas Oficiales Mexicanas (NOM) en M\'exico
        \item Emitidos en grandes vol\'umenes por m\'ultiples fuentes
        \item Riesgo regional para salud p\'ublica y ecosistemas
        \item Incluyen: 
        \begin{itemize}
            \item Di\'oxido de azufre (\ce{SO2})
            \item Di\'oxido de nitr\'ogeno (\ce{NO2})
            \item Mon\'oxido de carbono (\ce{CO})
            \item Part\'{\i}culas (PM$_{2.5}$/PM$_{10}$)
            \item Ozono troposf\'erico (\ce{O3})
            \item Plomo (\ce{Pb})
        \end{itemize}
    \end{itemize}
    
    \item \textbf{T\'oxicos ambientales}:
    \begin{itemize}
        \item Listado en el \emph{Clean Air Act} (189 compuestos) y en M\'exico es la  \gls{NOM}-165-\gls{SEMARNAT}-2013\footnote{Es el documento regulatorio que establece la lista de sustancias sujetas a reporte para el Registro de Emisiones y Transferencia de Contaminantes (RETC) en M\'exico. Esta norma, publicada el 24 de enero de 2014, tiene como objetivo regular las emisiones y transferencias de contaminantes al aire, agua, suelo y subsuelo, as\'{\i} como en residuos peligrosos y descargas de agua residual.}
        \item Efectos espec\'{\i}ficos: carcinog\'enesis, mutagenicidad, toxicidad reproductiva
        \item Ejemplos detallados en el \textbf{Cuadro~\ref{airtox}}
        \item Metodolog\'{\i}as de an\'alisis de riesgo estandarizadas de \gls{EPA} IRIS y  \gls{CARB} Hot Spots (AB 2588)
    \end{itemize}
\end{itemize}

Los m\'etodos cuantitativos de evaluaci\'on de riesgos aplican a toda sustancia t\'oxica atmosf\'erica, aunque los casos de estudio num\'ericos en este documento se enfocan espec\'{\i}ficamente en los compuestos listados en el \textbf{Cuadro~\ref{airtox}}.

 \begin{table}[!htbp]
\caption{Contaminantes T\'oxicos}
\begin{scriptsize}
\tablefirsthead{  \hline {\bf CAS} & {\bf Compuesto} &{\bf  CAS} &{\bf Compuesto}\\}
\begin{supertabular}{| p{0.10\textwidth}p{0.35\textwidth}p{0.10\textwidth}p{0.35\textwidth} |}\hline%
75070 & Acetaldeh\'{\i}do &133904 & Cloranfen \\
60355 & Acetamina  & 57749 & Clordano\\
75058 & Acetonitrilo & 7782505 & Cloro \\
98862 & Acetofenona & 79118 & Ac. clorac\'etico\\
53963 & 2-Acetilaminofluoreno & 532274 & 2-Cloracetafenona\\
107028 & Acroleina & 108907 & Clorobenceno\\
79061 & Acrilamina & 510156 & Clorobencilato\\
79107 & Ac. acr\'{\i}lico & 67663 & Cloroformo\\
107131 & Acrilonitrilo & 107302 & Clorometil metil eter\\
107051 & Alil cloruro & 126998 & Cloropreno \\
92671 & 4-aminobifenil & 1319773 & Cresoles (isomeros y mezclas)\\
62533 & Anilina &9587 &\emph{o}-Cresol\\
90040 & \emph{o}-Anisidina & 108394 &\emph{m}-Cresol \\
1332214 & Asbestos & 106445 &\emph{p}-Cresol \\
71432 & Benceno (incluyendo& 98828 &Cumeno \\
            & benceno de gasolina) & 94757 &2,4-diclorofenoxiac\'etico,  \\
92875 & Bencidina &  &sales y esteres\\
98077 & Benzotricloro & 3547044 & DDE\\
100447 & Cloruro de Bencilo & 334883 &  Diazometano\\
92524 & Bifenil & 132649 & Dibenzofuranos\\
117817 & Bis(2-etilhexil) ftalato (DEHP) & 96128 & 1,2-Dibromo- 3- cloropropano\\
542881& Bis(clorometil) eter &8742&Dibutilftalato\\
75252 & Bromformo &106467&1,4- Dicloro benceno (p) \\
106990 &1,3-Butadieno &91941& 3,3-Diclorobencidine \\
156627 & Cianamida de calcio &111444&Dicloroetil eter (bis(2-cloroetil)eter)\\
105602& Caprolactama& 542756 &1,3-Dicloropropeno\\ 
133062 & Captan & 62737 & Diclorvos (DDVP)\\
63252 & Carbaril &111422& Dietanolamina\\
75150 & Disulfuro de carbono & 121697 & N,N-Dietilanilina (NN-Dimetilanilina) \\
56235 & Tetracloruro de Carbono &64675&Sulfato de dietilo\\
463581 & Sulfuro de Carbonilo  &110543&Hexano\\
120809 & Catecol & 302012&Hidrazina \\
119904 & 3,3-Dimetoxybencidina &7647010&  \'acido Hidrocl\'orico\\
60117 & Dimetil ainoazoenceno&7664393& Fluoruro de hidr\'ogeno\\
119937 & 1,3-Dimetil bencidine & &(\'acido fluorh\'{\i}drico)\\
79447 & Clouro de Dimetil cabamol&123319&Hidroquinona\\
  68122 & Dimetil formamida &78591&Isoforona\\
  57147 & 1,2-Dimetil hidracina & 58899&Lindano ( todos sus is\'omeros) \\
  131113&  Dimetil ftalato &108316&Anh\'{\i}drico mal\'eico \\
  77781 &Dimetil sulfato &67561&Metanol \\
  534521& 4,6-Dinitro-o-cresol y sales & 72435&Metoxiclor \\
  51285 & 2,4- Dinitrofenol & 74839 & Bromuro de metilo (bromometano)\\
  121142 & 2,4- Dinitrotolueno & 74873& Cloruro de metilo \\
  123911 & 1,4 Dioxano (1,4-oxido de dietileno) &71556&Metil clororformo (1,1,1-tricloroeteno)\\
  122667 & 1,2-Difenilhidrazina &78933&Metil etil cetona (2-butanona) \\
  106898 & Epicloridina (1-cloro 2,3-epoxipropano)&60344 &Metil hidrazina\\
  106887 & 1,2-Epoxibutano &74884&Ioduro de metilo (iodometano)\\
  140885 & Etil acrilato &51796 &Carbamoato de etilo (uretano) \\
  100414 & Etil benceno &  &\\\hline
\end{supertabular}
\end{scriptsize}
\label{airtox}
\end{table}%


\subsection{Efectos a la salud}
Los efectos a la salud relacionados con la exposici\'on a los compuestos t\'oxicos pueden variar ampliamente. Un efecto espec\'{\i}fico en la salud usualmente se relaciona a un intervalo de niveles de exposici\'on. El \textbf{Cuadro~\ref{efectos}} lista algunos contaminantes y sus efectos asociados con la salud.
\index{contaminacion@contaminaci\'on!efectos a la salud}
 \begin{table}[htbp]
\caption{Efectos a la salud de algunos contaminantes}
\begin{center}
{\small
\begin{tabular}{|l|l|}\hline
{\bf Contaminante} &{\bf  Efectos a la salud}\\\hline
Ozono   & Problemas en el tracto respiratorio como la dificultad de respirar,  \\
               &  reducci\'onde la funci\'on respiratoria, a la resistencia a la infecci\'on, \\
               & asma, irritaci\'on de ojos, congesti\'on,  y posible envejecimiento \\
               & del tejido pulmonar.\\ \index{ozono!efectos a la salud}
Part\'{\i}culas  & Irritaci\'on de ojos y garganta, bronquitis, da\~no al pulm\'on\\
               & reducci\'on de la visibilidad.\\ \index{material particulado} \index{particulas@part\'{\i}culas!efectos a la salud}
CO        & Incapacidad de la sangre para transportar ox\'{\i}geno; afectaci\'on\\
              & a los sistemas cardiovascular, nervioso y pulmonar\\ \index{monoxido@mon\'oxido!de carbono!efectos a la salud}
\ce{SO2}& Problemas en el tracto respiratorio, da\~no permanente a los\\
             & tejidos del pulm\'on\\  \index{SO2@\ce{SO2}!efectos a la salud}   
Plomo  & Retardo y da\~no cerebral, especialmente en ni\~nos\\\index{plomo}
\ce{NO2}& Enfermedades respiratorias y da\~no al pulm\'on\\ \index{no2@\ce{NO2}!efectos a la salud}
Asbesto & Variedad de enfermedades de pulm\'on, particularmente c\'ancer \\\index{asbesto!efectos a la salud}
Berilio  & Primeramente da\~no al pulm\'on, aunque tambi\'en da\~na al h\'{\i}gado,\\ \index{berilio!efectos a la salud}
             &  bazo, ri\~nones y gl\'andulas linf\'aticas.\\ \index{higado@h\'{\i}gado}
Mercurio& Varias \'areas del cerebro como tambi\'en hay afectaci\'on en los \\
                & ri\~nones e intestinos\\ \index{mercurio!efectos a la salud}
Cloruro de vinilo& C\'ancer de pulm\'on e h\'{\i}gado\\ \index{vinilo! cloruro de}
Ars\'enico & C\'ancer\\ \index{arsenico@ars\'enico!efectos a la salud}
Radion\'ucleos& C\'ancer\\ \index{radionucleos@radion\'ucleos!efectos a la salud}
Benceno & Leucemia\\ \index{benceno!efectos a la salud}
Emisiones de coke & C\'ancer en vias respiratorias\\ \hline
\end{tabular}}
\end{center}
\label{efectos}
\end{table}%

\clearpage
% =============================================================================
% exercises/ejercicios_cap06.tex
% Ejercicios propuestos — Capítulo 6: Contaminación Atmosférica
% =============================================================================

\section*{Ejercicios Propuestos}
\addcontentsline{toc}{section}{Ejercicios propuestos}
\label{sec:ejercicios_cap06}

\begin{bloque_ejercicios}[Ejercicios --- Cap\'{\i}tulo 5: Contaminaci\'on Atmosf\'erica]

\begin{ejercicios}

% ----- Ejercicio 1: Fuentes y clasificación de contaminantes ----------------
\item Clasifique cada uno de los siguientes contaminantes indicando:
(i)~su \textbf{origen} (biol\'ogico, tierra s\'olida, oc\'eano o formaci\'on
in situ), (ii)~si es \textbf{primario o secundario}, y
(iii)~si pertenece a los \textbf{contaminantes criterio} o no:

{%
\centering
\begin{tabular}{clccc}
\toprule
\textbf{N\'um.} & \textbf{Contaminante} & \textbf{Origen} &
\textbf{Primario/Secundario} & \textbf{Criterio} \\
\midrule
a) & \ce{SO2} (erupci\'on volc\'anica) & & & \\
b) & \ce{O3} troposf\'erico        & & & \\
c) & \ce{CH4} (humedales)          & & & \\
d) & \ce{NO} (combusti\'on)         & & & \\
e) & \ce{PAN}                       & & & \\
f) & \ce{H2SO4} (aerosol)           & & & \\
g) & \ce{Pb} (gasolina)             & & & \\
h) & DMS (fitoplancton)             & & & \\
\bottomrule
\end{tabular}
\par
}%

% ----- Ejercicio 2: Modelo de difusión de COS -------------------------------
\item Con base en el modelo de intercambio de COS presentado en el
cap\'{\i}tulo, en el que la concentraci\'on decrece exponencialmente con la
profundidad $z$ seg\'un $C(z)=C(0)\exp(-z/H)$, con altura de escala
$H=\left(D/k\right)^{1/2}$:

\begin{enumerate}[label=\alph*)]
    \item Si el coeficiente de difusi\'on turbulenta se duplica a
          $D = \SI{0.010}{\metre\squared\per\second}$ y la constante de
          destrucci\'on permanece en $k = 5.0\times10^{-6}\,\si{\per\second}$,
          calcule la nueva altura de escala $H'$ y la profundidad a la que
          $C(z) = C(0)/2$.

    \item Demuestre que la vida media de una mol\'ecula de COS en el oc\'eano
          es siempre $\tau = 1/k$, independientemente del valor de $D$.
          \textit{(Siga el procedimiento del ejemplo del cap\'{\i}tulo.)}

    \item ?`C\'omo afectar\'{\i}a una disminuci\'on de la productividad biol\'ogica
          marina al flujo de COS hacia la atm\'osfera? Argumente
          cualitativamente en t\'erminos de $C(0)$ y la densidad de columna.
\end{enumerate}

\textit{[Respuesta parcial: (a) $H' = \SI{44.7}{\metre}$;
profundidad para $C(0)/2$: $z = H'\ln 2 \approx \SI{31.0}{\metre}$]}

% ----- Ejercicio 3: Ciclo geoquímico del CO2 --------------------------------
\item El ciclo de meteorizaci\'on de la calcita (\ce{CaCO3}) involucra la
reacci\'on neta:
\[
    \ce{CaCO3 (s) + CO2 (g) + H2O (l) <=> Ca^{2+} (ac) + 2HCO3^- (ac)}
\]

\begin{enumerate}[label=\alph*)]
    \item Identifique en qu\'e direcci\'on procede la reacci\'on
          (derecha o izquierda) durante: (i) la meteorizaci\'on continental
          por lluvias \'acidas, y (ii)~la precipitaci\'on de sedimentos
          marinos. Justifique.

    \item Cuando el \ce{CO2} se disuelve en agua pura se forman tres
          especies en equilibrio. Escriba las tres ecuaciones de equilibrio
          en secuencia y la reacci\'on neta, tal como se presenta en el
          ejemplo del cap\'{\i}tulo.

    \item El tiempo de residencia del \ce{HCO3^-} en los oc\'eanos es
          $\tau \approx 7.5\times10^4$~a\~nos. Si la tasa de entrada de
          \ce{HCO3^-} al oc\'eano es $F = \SI{3.0e12}{\mole}/$~a\~no,
          estime la cantidad total de \ce{HCO3^-} en el oc\'eano usando
          $M = F \cdot \tau$.

    \item Discuta cu\'al ser\'{\i}a el efecto de un incremento sostenido de
          \ce{CO2} atmosf\'erico (debido a actividades antropog\'enicas)
          sobre el equilibrio de la reacci\'on neta anterior y sobre
          la acidez de los oc\'eanos.
\end{enumerate}

\textit{[Respuesta parcial:
(c) $M = 3.0\times10^{12} \times 7.5\times10^{4}
    = 2.25\times10^{17}\,\si{\mole}$]}

% ----- Ejercicio 4: Carboxihemoglobina (CO) ---------------------------------
\item La concentraci\'on de carboxihemoglobina (\%COHb) en sangre se
estima con la ecuaci\'on:
\[
    \%\ce{COHb} = 0.005\,[\ce{CO}]^{0.85}\,(\alpha t)^{0.63}
\]
donde $[\ce{CO}]$ est\'a en ppm, $\alpha$ es el coeficiente de actividad
y $t$ es el tiempo de exposici\'on en minutos.

\begin{enumerate}[label=\alph*)]
    \item Una persona sedentaria ($\alpha = 0.94$) es expuesta a
          $[\ce{CO}] = 35$~ppm durante 2 horas.
          Calcule \%COHb e indique qu\'e efectos fisiol\'ogicos
          se esperan seg\'un los umbrales del cap\'{\i}tulo.

    \item Un trabajador de construcci\'on ($\alpha = 3.0$) labora en un
          espacio confinado con $[\ce{CO}] = 50$~ppm durante 90~min.
          Calcule \%COHb y eval\'ue el riesgo.

    \item Usando los resultados de los incisos anteriores, explique por
          qu\'e el mismo nivel de \ce{CO} resulta m\'as peligroso para
          personas con actividad f\'{\i}sica intensa que para personas
          en reposo.

    \item Estacionamiento subterr\'aneo. La norma mexicana NOM-021 fija
          el l\'{\i}mite de \%COHb en trabajadores en $\leq 5\%$.
          ?`Cu\'al es la concentraci\'on m\'axima permisible de \ce{CO} (en ppm)
          para un trabajador sedentario ($\alpha = 0.94$) expuesto 8 horas?
          \textit{Despeje $[\ce{CO}]$ de la ecuaci\'on.}
\end{enumerate}

\textit{[Respuesta parcial:
(a) $\%\text{COHb} \approx 3.1\%$, efectos ligeros;
(b) $\%\text{COHb} \approx 7.5\%$, efectos psicomotores;
(d) $[\ce{CO}]_{\max} \approx 17$~ppm]}

% ----- Ejercicio 5: Compuestos de azufre y lluvia ácida ---------------------
\item El di\'oxido de azufre emitido a la atm\'osfera se oxida
progresivamente y puede generar lluvia \'acida.

\begin{enumerate}[label=\alph*)]
    \item Escriba la secuencia de reacciones por las que el \ce{SO2}
          atmosf\'erico se convierte en \'acido sulf\'urico (\ce{H2SO4})
          en presencia de agua y ox\'{\i}geno. Indique en cu\'al paso
          interviene el \ce{OH} radical.

    \item Una central termoelect\'ectrica quema carb\'on con contenido
          de azufre del 2.5\% en masa. Si la tasa de consumo es de
          \SI{500}{\mega\gram\per\hour}, estime la emisi\'on de \ce{SO2}
          en \si{\mega\gram\per\hour}. Dato: masa molar \ce{S} = 32~g/mol;
          \ce{SO2} = 64~g/mol.

    \item Seg\'un los umbrales de efectos en la salud del cap\'{\i}tulo,
          clasifique el riesgo para la poblaci\'on circundante si
          la concentraci\'on ambiente de \ce{SO2} alcanza 0.5~ppm
          durante 3 horas.

    \item Explique c\'omo el \ce{SO2} contribuye al incremento de
          PM$_{10}$ y PM$_{2.5}$, identificando la especie qu\'{\i}mica
          intermedia formada.
\end{enumerate}

\textit{[Respuesta parcial:
(b) Masa de S quemado = $500 \times 0.025 = \SI{12.5}{\mega\gram\per\hour}$;
emisi\'on de \ce{SO2} $= 12.5 \times (64/32) = \SI{25}{\mega\gram\per\hour}$;
(c) $0.5 > 0.3$~ppm $\Rightarrow$ riesgo alto: enfermedades
cardiorrespiratorias]}

% ----- Ejercicio 6: Ozono troposférico y estratosférico ---------------------
\item Con base en el \textbf{Cuadro~\ref{ozon:1}} y el texto del cap\'{\i}tulo:

\begin{enumerate}[label=\alph*)]
    \item Complete la siguiente tabla comparativa indicando
          \textit{altitud}, \textit{porcentaje del ozono total},
          \textit{efecto principal} y \textit{problem\'atica actual}
          para cada tipo de ozono:

{%
\centering
\begin{tabular}{lcc}
\toprule
\textbf{Caracter\'{\i}stica} & \textbf{Troposf\'erico} & \textbf{Estratosf\'erico}\\
\midrule
Altitud (km)         & & \\
\% del total         & & \\
Efecto principal     & & \\
Problem\'atica actual & & \\
\bottomrule
\end{tabular}
\par
}%

    \item La formaci\'on de ozono troposf\'erico requiere de COV,
          NO$_x$ y luz solar. Explique por qu\'e en una ciudad con
          alta emisi\'on de NO pero baja de COV la concentraci\'on de
          \ce{O3} puede ser menor que en los suburbios.

    \item Mencione dos factores meteorol\'ogicos u orogr\'aficos que
          favorecen la acumulaci\'on de ozono troposf\'erico en
          ciudades como la Ciudad de M\'exico o Los \'Angeles,
          con base en la informaci\'on del cap\'{\i}tulo.

    \item El ozono estratosf\'erico absorbe la radiaci\'on UV.
          ?`Qu\'e consecuencias para la salud humana se derivar\'{\i}an
          de una reducci\'on significativa de esta capa?
          Mencione al menos dos efectos.
\end{enumerate}

% ----- Ejercicio 7: Partículas y visibilidad --------------------------------
\item El alcance \'optico $V$ (visibilidad horizontal) se relaciona
con el coeficiente de extinci\'on $\beta_{ext}$ mediante la f\'ormula de
Koschmieder:
\[
    V = \frac{3.912}{\beta_{ext}} \quad [\si{\kilo\metre}]
\]

\begin{enumerate}[label=\alph*)]
    \item Calcule la visibilidad en un d\'{\i}a pristino con
          $\beta_{ext} = 10\ \text{Mm}^{-1}$ y en un episodio de
          contaminaci\'on severa con
          $\beta_{ext} = 1{,}000\ \text{Mm}^{-1}$.

    \item Seg\'un el cap\'{\i}tulo, la humedad relativa mayor a 70\%
          incrementa $\beta_{ext}$ entre 2 y 5 veces por crecimiento
          higrosc\'opico de las part\'{\i}culas. Si en una ma\~nana con
          HR~$= 85\%$ el $\beta_{ext}$ seco es de $200\ \text{Mm}^{-1}$,
          estime el intervalo de visibilidad esperado.

    \item Las PM$_{2.5}$ tienen di\'ametros de \SIrange{0.12}{0.22}{\micro\metre},
          que coinciden con la condici\'on $\alpha \approx 1$ para luz
          visible ($\lambda \approx \SI{0.55}{\micro\metre}$).
          Explique en t\'erminos del par\'ametro de Mie
          $\alpha = 2\pi r/\lambda$ por qu\'e estas part\'{\i}culas
          dispersan la luz de manera m\'axima.

    \item Clasifique los tres gases (\ce{SO2}, \ce{NO}, \ce{H2O})
          que menciona el cap\'{\i}tulo como influyentes en la visibilidad,
          describiendo el mecanismo espec\'{\i}fico de cada uno.
\end{enumerate}

\textit{[Respuesta parcial:
(a) $V_{pristino} = 3.912/10 \times 10^{-3}
    = \SI{391.2}{\kilo\metre}$;
    $V_{contam} = \SI{3.9}{\kilo\metre}$;
(b) $\beta_{ext}^{humedo} \in [400, 1000]\ \text{Mm}^{-1}$;
    $V \in [3.9, 9.8]\ \si{\kilo\metre}$]}

% ----- Ejercicio 8: Plomo en la atmósfera -----------------------------------
\item El tetraetilo de plomo (\ce{(C2H5)4Pb}) se a\~nad\'{\i}a a la gasolina
a raz\'on de \SI{1.1}{\gram}/gal (equivalente a\SI{0.29}{\gram\per\litre}).

\begin{enumerate}[label=\alph*)]
    \item Si un veh\'{\i}culo consume \SI{40}{\litre}/semana de gasolina
          con esta formulaci\'on, calcule la masa de plomo emitida
          a la atm\'osfera por semana. Dato: fracci\'on m\'asica de Pb en
          \ce{(C2H5)4Pb} $= 207/(207 + 4 \times 29) = 64.1\%$.

    \item Un ni\~no de 5 a\~nos presenta concentraci\'on sangu\'{\i}nea de
          plomo de \SI{15}{\micro\gram\per\deci\liter}. Seg\'un los umbrales
          del cap\'{\i}tulo (CDC 2023), ?`este valor est\'a por debajo o
          por encima del umbral de da\~nos hematol\'ogicos?
          ?`Qu\'e efectos neurol\'ogicos se podr\'{\i}an observar?

    \item Explique el mecanismo por el cual el pH del agua potable
          influye en la concentraci\'on de plomo ingerido a partir
          de tuber\'{\i}as de plomo. ?`Un agua con pH~$= 6$ es m\'as o
          menos riesgosa que una con pH~$= 8$?

    \item Desde el punto de vista de la pol\'{\i}tica ambiental, ?`por
          qu\'e la eliminaci\'on del plomo en la gasolina se considera
          una de las medidas de salud p\'ublica de mayor impacto en
          el siglo XX? Argumente con base en las v\'{\i}as de exposici\'on
          descritas en el cap\'{\i}tulo.
\end{enumerate}

\textit{[Respuesta parcial:
(a) Masa TEL = $40 \times 0.29 = \SI{11.6}{\gram}/$semana;
masa Pb $= 11.6 \times 0.641 = \SI{7.4}{\gram}/$semana;
(b) \SI{15}{\micro\gram\per\deci\liter} $<$ \SI{20}{\micro\gram\per\deci\liter}: por debajo de umbral
hematol\'ogico, pero en zona de riesgo neurol\'ogico]}

% ----- Ejercicio 9: Esmog y contaminantes primarios/secundarios -------------
\item El esmog fotoqu\'{\i}mico se genera cuando contaminantes primarios
reaccionan bajo la influencia de la luz solar.

\begin{enumerate}[label=\alph*)]
    \item Describa las diferencias fundamentales entre el
          \textbf{esmog industrial} (tipo Londres) y el
          \textbf{esmog fotoqu\'{\i}mico} (tipo Los \'Angeles) en cuanto a:
          condiciones meteorol\'ogicas, contaminantes involucrados
          y efectos en la salud.

    \item La reacci\'on simplificada de formaci\'on de esmog
          fotoqu\'{\i}mico puede escribirse como:
          \[
            \text{COV} + \text{NO}_x + h\nu \longrightarrow
            \text{esmog fotoqu\'{\i}mico} (\ce{O3},\,\text{PAN},\ldots)
          \]
          Indique qu\'e es el PAN\index{PAN}, a qu\'e clase de contaminante
          pertenece (primario o secundario) y por qu\'e resulta
          especialmente irritante.

    \item En episodios de esmog, las concentraciones de \ce{NO}
          son m\'aximas en la ma\~nana y decaen al mediod\'{\i}a,
          mientras que el \ce{O3} presenta el comportamiento opuesto.
          Explique este ciclo diurno en t\'erminos de las reacciones
          del cap\'{\i}tulo.

    \item ?`Por qu\'e las ciudades con cuencas rodeadas de monta\~nas
          (como la Ciudad de M\'exico) son especialmente vulnerables
          a los episodios de esmog fotoqu\'{\i}mico?
\end{enumerate}

% ----- Ejercicio 10 (avanzado): Compuestos tóxicos y análisis de riesgo ----
\item \textbf{[Avanzado]} El benceno (\ce{C6H6}) es uno de los compuestos
t\'oxicos listados en el \textbf{Cuadro~\ref{airtox}} y es un agente
canser\'ogeno reconocido.

\begin{enumerate}[label=\alph*)]
    \item Clasifique al benceno seg\'un las dos categor\'{\i}as de
          contaminantes descritas en la Secci\'on~\ref{ctox}
          (\textit{contaminante criterio} vs.\ \textit{t\'oxico ambiental})
          e indique el efecto espec\'{\i}fico sobre la salud reportado en
          el \textbf{Cuadro~\ref{efectos}}.

    \item El benceno proviene principalmente de la evaporaci\'on
          de gasolina y de la combusti\'on incompleta. Identifique
          en qu\'e secci\'on del cap\'{\i}tulo se menciona la quema de
          biomasa como fuente de \ce{CH3Cl} y relacione ese proceso
          con la emisi\'on de benceno: ?`ambos son fuentes
          biol\'ogicas, antropog\'enicas o mixtas?

    \item La metodolog\'{\i}a EPA~IRIS establece una concentraci\'on
          de riesgo unitario (URF) para benceno de
          $\text{URF} = 7.8\times10^{-6}\,(\mu\text{g/m}^3)^{-1}$.
          El riesgo individual de c\'ancer se calcula como:
          \[
              R = \text{URF} \times C \times \frac{ED}{70}
          \]
          donde $C$ es la concentraci\'on en $\mu$g/m$^3$,
          $ED$ es la duraci\'on de la exposici\'on en a\~nos y 70 es
          la expectativa de vida. Calcule $R$ para una persona que
          vive 30 a\~nos en una zona con $C = \SI{5}{\micro\gram\per\cubic\metre}$.

    \item Si la EPA considera ``aceptable'' un riesgo individual
          $R \leq 1\times10^{-6}$ (1 caso en un mill\'on),
          ?`el valor calculado en el inciso anterior es aceptable?
          ?`Qu\'e concentraci\'on m\'axima de benceno ser\'{\i}a permisible
          para cumplir este criterio en las mismas condiciones
          de exposici\'on?
\end{enumerate}

\textit{[Respuesta parcial:
(c) $R = 7.8\times10^{-6} \times 5 \times (30/70)
    \approx 1.67\times10^{-5}$;
(d) No aceptable; $C_{\max} = 10^{-6}/(7.8\times10^{-6}\times 30/70)
    \approx \SI{0.30}{\micro\gram\per\cubic\metre}$]}

\end{ejercicios}

\end{bloque_ejercicios}
% =============================================================================
% FIN DE ejercicios_cap06.tex
% =============================================================================

 
